\section{The collected total in one function basis}
\label{sec:comp-total-energy-collected}
All contributions can now be combined in the same energy variable. This section gives the complete coefficient of $\alpha_{\mathrm T}^3/m^2$, including the virtual correction and the integrated two-photon final state. It is obtained from the independently evaluated sectors above. Define $X=s_{\mathrm{phys}}/m^2-1>0$ and $Y=\sqrt{X/(X+4)}$. Use the GPL kernels and real branches in Section~\ref{sec:comp-virtual-energy-basis}, including the removable combinations containing the letter $1$. The complete result is
\begin{equation}
\sigma_{\mathrm{inc}}^{\mathrm{NLO}}=\frac{\pi\alpha_{\mathrm T}^2}{m^2}B(X)+\frac{\alpha_{\mathrm T}^3}{m^2}C_1(X)+\frac1{m^2}O(\alpha_{\mathrm T}^4),
\end{equation}
where $B$ is given in \eqref{eq:comp-simplified-born} and
\begin{equation}\boxed{C_1(X)=\sum_{j=0}^{22}c_j(X,Y)F_j(X).}\label{eq:comp-total-energy-collected}\end{equation}
The functions and their full coefficients follow. No unevaluated phase-space or Feynman-parameter integral occurs in this expression.
\begin{equation}
 F_{0}=1,\qquad c_{0}=\frac{- 273 X^{3} + 982 X^{2} + 2960 X + 1744}{24 X^{2} \left(X + 1\right)^{2}}.
\end{equation}
\begin{equation}
 F_{1}=G(-1;X),\qquad c_{1}=\frac{37 X^{4} - 54 X^{3} - 339 X^{2} - 428 X - 184}{4 X^{3} \left(X + 1\right)^{2}}.
\end{equation}
\begin{equation}
 F_{2}=G(0;X),\qquad c_{2}=\frac{- 14 X^{4} - 17 X^{3} + 17 X^{2} + 22 X + 8}{2 X \left(X - 1\right) \left(X + 1\right)^{3}}.
\end{equation}
\begin{equation}
 F_{3}=G(-1,-1;X),\qquad c_{3}=\frac{\begin{aligned}&- 4 X^{6} - 35 X^{5} + 31 X^{4} + 755 X^{3}\\&+1765 X^{2} + 1506 X + 440\end{aligned}}{X^{3} \left(X + 1\right)^{2} \left(X + 4\right)}.
\end{equation}
\begin{equation}
 F_{4}=G(-1,0;X),\qquad c_{4}=\frac{X^{4} + 7 X^{3} + X^{2} - 3 X - 2}{X^{3} \left(X + 1\right)^{2}}.
\end{equation}
\begin{equation}
 F_{5}=G(-2,-1;X),\qquad c_{5}=\frac{\begin{aligned}&- 4 X^{5} - 104 X^{4} - 584 X^{3} - 1264 X^{2}\\&- 1152 X - 384\end{aligned}}{X^{3} \left(X + 1\right)^{2} \left(X + 4\right)}.
\end{equation}
\begin{equation}
 F_{6}=G(0,-1;X),\qquad c_{6}=\frac{\begin{aligned}&X^{6} + 8 X^{5} - 17 X^{4} - 210 X^{3}\\&- 448 X^{2} - 352 X - 96\end{aligned}}{X^{3} \left(X + 1\right)^{2} \left(X + 4\right)}.
\end{equation}
\begin{equation}
 F_{7}=G(1,0;X),\qquad c_{7}=\frac{- X^{2} + X - 2}{X^{3}}.
\end{equation}
\begin{equation}
 F_{8}=G(r,-1;X),\qquad c_{8}=\frac{Y \left(3 X^{4} + 18 X^{3} + 44 X^{2} - 8 X - 64\right)}{X^{4}}.
\end{equation}
\begin{equation}
 F_{9}=G(-1,-1,-1;X),\qquad c_{9}=\frac{- 2 X^{2} + 32 X + 46}{X^{3}}.
\end{equation}
\begin{equation}
 F_{10}=G(-1,-1,0;X),\qquad c_{10}=\frac{\left(2 X + 3\right) \left(4 X - 8\right)}{X^{3}}.
\end{equation}
\begin{equation}
 F_{11}=G(-1,-2,-1;X),\qquad c_{11}=\frac{8 X^{2} - 32 X - 48}{X^{3}}.
\end{equation}
\begin{equation}
 F_{12}=G(-1,0,-1;X),\qquad c_{12}=\frac{4 X^{2} - 14 X - 24}{X^{3}}.
\end{equation}
\begin{equation}
 F_{13}=G(-4,r,-1;X),\qquad c_{13}=- \frac{6 Y \left(X^{3} - 12 X - 8\right)}{X^{4}}.
\end{equation}
\begin{equation}
 F_{14}=G(0,-1,-1;X),\qquad c_{14}=\frac{- 5 X^{2} - 32 X + 8}{X^{3}}.
\end{equation}
\begin{equation}
 F_{15}=G(0,-1,0;X),\qquad c_{15}=\frac{2 X^{2} + 4 X - 12}{X^{3}}.
\end{equation}
\begin{equation}
 F_{16}=G(0,1,0;X),\qquad c_{16}=\frac{2 X^{2} + 4 X - 12}{X^{3}}.
\end{equation}
\begin{equation}
 F_{17}=G(0,r,-1;X),\qquad c_{17}=- \frac{3 Y \left(X - 2\right) \left(X + 4\right)}{X^{3}}.
\end{equation}
\begin{equation}
 F_{18}=G(r,-1,-1;X),\qquad c_{18}=\frac{3 Y \left(5 X^{4} + 14 X^{3} - 96 X^{2} - 352 X - 128\right)}{X^{4} \left(X + 4\right)}.
\end{equation}
\begin{equation}
 F_{19}=G(r,-1,0;X),\qquad c_{19}=- \frac{6 Y \left(X - 2\right) \left(X + 4\right)}{X^{3}}.
\end{equation}
\begin{equation}
 F_{20}=G(r,-2,-1;X),\qquad c_{20}=- \frac{16 Y \left(X^{4} + 2 X^{3} - 24 X^{2} - 80 X - 48\right)}{X^{4} \left(X + 4\right)}.
\end{equation}
\begin{equation}
 F_{21}=G(r,0,-1;X),\qquad c_{21}=- \frac{8 Y \left(X^{4} + 3 X^{3} - 18 X^{2} - 68 X - 24\right)}{X^{4} \left(X + 4\right)}.
\end{equation}
\begin{equation}
 F_{22}=G(r,r,-1;X),\qquad c_{22}=\frac{9 X^{2} - 24}{X^{3}}.
\end{equation}
\paragraph{Collection of the real terms.} Put $L=\ln(1+X)$, $M=\ln X$, and $J=\ln(1+X/2)$. The classical functions in the four-sector real expression satisfy
\begin{align}
\operatorname{Li}_2\!\left(\frac1{1+X}\right)&=\zeta(2)+G(-1,0;X)-G(-1,-1;X),\\
\operatorname{Li}_2\!\left(\frac2{2+X}\right)+\ln2\,J&=\zeta(2)+G(-2,0;X)-G(-2,-2;X),\\
\zeta(3)-\operatorname{Li}_3\!\left(\frac1{1+X}\right)&=\zeta(2)G(-1;X)+G(-1,-1,0;X)-G(-1,-1,-1;X).
\end{align}
Each identity follows by differentiating with respect to $X$. The constants are fixed as $X\to0^+$, where the dilogarithms and trilogarithm take their values at one. Every nonconstant GPL on these right-hand sides vanishes in this limit. In particular, the logarithmic lower endpoint in $G(-1,0;X)$ and $G(-2,0;X)$ is integrable because $\int_0^X|\ln t|\,\mathrm{d}t=O(X[1+|\ln X|])$. The remaining products of $L$, $M$, and $J$ are expanded with the shuffle identity, counting each interleaving even when repeated letters give the same word. All classical arguments above lie in $(0,1)$ for $X>0$, so this conversion introduces no continuation across a cut.
\paragraph{Assembly.} Substitute these identities, the completed $UW$ expression, and \eqref{eq:comp-ck-energy-compact} into \eqref{eq:comp-simplified-master}. The finite virtual coefficient is its constant term in $\ell_\lambda=\ln(\lambda_{\mathrm{phys}}/m)$ after the previously proved infrared cancellation. Taking this coefficient does not set the photon mass to $m$. The finite conversion term remains $-A(1+X)[\ln X+\tfrac12\ln(1+X)]$. Expanding this product gives the last coefficients needed in the sum above.
There are five contributions to this collection: $\widehat C_{V+\mathrm{UV}}$, $F_{UW}^{\eta}/[4X(1+X)]$, $\mathcal R_4(1+X)/[8X(1+X)]$, $-\mathcal C_K(1+X)/(1+X)$, and the finite conversion term. The factor $1/4$ for $UW$ follows from $\mathcal R_{\mathrm{rem}}=\mathcal R_4+2F_{UW}^{\eta}$, so its Bose-related second copy is already included. Expand each contribution in the same words and add their rational coefficients. This is an algebraic reconstruction and does not require a claim that the functions are linearly independent. The infrared coefficient used in the last contribution is \eqref{eq:comp-ir-short}, with $L^2=2G(-1,-1;X)$ and $\mathcal G(X)=G(r,-1;X)$.
The coefficients of five words vanish identically in the complete sum. In particular the coefficient of $G(-1,r,-1;X)$ from the finite soft term cancels its coefficient from the infrared conversion term. All coefficients proportional to $\pi^2$ outside the GPLs also cancel. The remaining coefficients are the rational functions listed above.
\paragraph{The apparent point $X=1$.} The factor $1/(X-1)$ multiplies $G(0;X)=\ln X$, so it has a removable limit. The product involving this coefficient tends to $+1$. The two GPLs containing the rational letter $+1$ are also real and finite:
\begin{align}
G(1,0;X)&=\ln X\ln|1-X|+\operatorname{Re}\operatorname{Li}_2(X),\\
G(0,1,0;X)&=2\operatorname{Re}\operatorname{Li}_3(X)-\ln X\operatorname{Re}\operatorname{Li}_2(X).
\end{align}
Their derivatives are $\ln X/(X-1)$ and $G(1,0;X)/X$. At $X=1$ their values are $\zeta(2)$ and $2\zeta(3)$. Thus the apparent inner pole does not require a contour choice for the physical answer.
